\documentclass[11pt]{article}
\usepackage[UTF8, fontset=fandol, scheme=plain]{ctex}
\usepackage[margin=1in]{geometry}
\usepackage{amsmath,amssymb,amsthm}
\usepackage{array}
\usepackage{parskip}

\newtheorem{theorem}{Theorem}
\newtheorem{lemma}[theorem]{Lemma}
\newtheorem{corollary}[theorem]{Corollary}
\newtheorem{proposition}[theorem]{Proposition}
\theoremstyle{definition}
\newtheorem{definition}[theorem]{Definition}
\theoremstyle{remark}
\newtheorem{remark}[theorem]{Remark}

\newcommand{\Cfive}{C_5}

\title{Disproof of conjecture \texttt{00000003485}}
\author{earthking11}
\date{2026-09-13}

\begin{document}
\sloppy
\maketitle

\section*{Verdict: FALSE}

We disprove conjecture \texttt{00000003485} as stated. The refutation is
unconditional and elementary, and the witness is the smallest possible: the
directed $5$-cycle. It has out-degree $1$ everywhere, directed girth $5$, and
yet no kernel. The same failure occurs for every directed odd cycle, and the
failure is stable under strongly connected blow-ups, so it is not an artefact of
out-degree $1$.

\section{The conjecture}

Quoted verbatim from \texttt{conjectures/00000003485.md}:

\begin{quote}
\textbf{Definition:} A kernel of a digraph is an independent out-stable set,
with the existence criterion the core problem. \textbf{Conjecture:} Digraphs
with bounded out-degree and girth at least five have kernels; the critical
configuration of the condition is the directed odd cycle with a one-point
source, and the bound is a tightening of Richardson's theorem. (kernel existence
strengthened condition)
\end{quote}

\begin{quote}
\textbf{定义：}有向图的核为独立外稳定集，其存在性判据是核心问题。\textbf{猜想：}出度有界且围长至少五的有向图存在核；条件的临界构型为有向奇圈加单点源且该界为
Richardson 定理的紧化。（核存在强化条件）
\end{quote}

\begin{definition}[digraph]
A digraph $D = (V, A)$ consists of a vertex set $V$ and an arc set
$A \subseteq V \times V$.
\end{definition}

\begin{definition}[kernel]
A \emph{kernel} of $D = (V, A)$ is a set $S \subseteq V$ that is
\begin{enumerate}
\item \emph{independent}: no arc joins two members of $S$ in either direction,
      i.e.\ for all $u, v \in S$ neither $(u, v) \in A$ nor $(v, u) \in A$; and
\item \emph{out-stable} (absorbing): every vertex outside $S$ has an arc into
      $S$, i.e.\ for every $v \in V \setminus S$ there is $s \in S$ with
      $(v, s) \in A$.
\end{enumerate}
\end{definition}

\begin{definition}[out-degree, girth]
The \emph{out-degree} of $v$ is $d^+(v) = |\{w : (v,w) \in A\}|$. The
\emph{directed girth} is the length of the shortest directed cycle, i.e.\ the
least $k \ge 1$ with a closed directed walk of length $k$; it is $+\infty$ for
an acyclic digraph. The \emph{(underlying undirected) girth} is the length of
the shortest cycle in the underlying undirected graph.
\end{definition}

The hypothesis ``bounded out-degree'' is read in its weakest useful form, as a
global bound $d^+(v) \le B$ for a fixed constant $B$. A universal implication of
the form ``every digraph with bounded out-degree and girth $\ge 5$ has a
kernel'' is refuted by a single digraph satisfying both hypotheses and having no
kernel.

\section{The witness: the directed 5-cycle}

\begin{definition}[$\Cfive$]
Let $\Cfive$ be the digraph with vertex set $V = \{0, 1, 2, 3, 4\}$ and arc set
\[
A = \{\, (i,\; i + 1 \bmod 5) \;:\; i \in V \,\}.
\]
It is the directed cycle $0 \to 1 \to 2 \to 3 \to 4 \to 0$.
\end{definition}

\begin{lemma}[hypotheses hold]\label{lem:hyp}
The digraph $\Cfive$ has out-degree $1$ at every vertex, hence bounded
out-degree, and both its directed girth and its underlying undirected girth are
equal to $5 \ge 5$.
\end{lemma}

\begin{proof}
Every vertex has exactly one out-neighbour, its cyclic successor, so
$d^+(v) = 1$ for all $v$; in particular the out-degree is bounded (take
$B = 1$). The directed closed walk $0 \to 1 \to 2 \to 3 \to 4 \to 0$ has length
$5$, so the directed girth is at most $5$. There is no closed directed walk of
length $1, 2, 3$, or $4$: every arc moves the index up by exactly $1$ modulo
$5$, so a closed walk of length $k$ would force $k \equiv 0 \pmod 5$, and no
positive $k < 5$ satisfies this. Hence the directed girth is exactly $5$. The
underlying undirected graph is the ordinary cycle $0$--$1$--$2$--$3$--$4$--$0$,
whose girth is $5$.
\end{proof}

\section{$\Cfive$ has no kernel}

\begin{theorem}\label{thm:no-kernel}
The directed $5$-cycle $\Cfive$ has no kernel.
\end{theorem}

\begin{proof}
We use two elementary facts about independence and absorption in $\Cfive$.

\emph{Independence.} An independent set contains no two cyclically consecutive
vertices, because $(i, i+1) \in A$ for every $i$. Hence an independent set is an
independent set of the underlying undirected $5$-cycle, whose independence
number is $2$; so $|S| \le 2$, and if $|S| = 2$ then $S = \{i, i+2\}$ for some
$i$ (the only non-adjacent pairs are the pairs at cyclic distance $2$).

\emph{Absorption.} If $v \notin S$, the unique arc leaving $v$ is
$(v, v+1)$, so $v$ can have an arc into $S$ only if $v + 1 \in S$. Thus
out-stability reads:
\begin{equation}\label{eq:absorb}
\text{for every } v \in V: \quad v \in S \;\text{ or }\; v + 1 \in S .
\end{equation}

Now let $S$ be independent, and consider its size.

\begin{itemize}
\item $|S| = 0$: then $0 \notin S$ and $0 + 1 = 1 \notin S$, violating
      \eqref{eq:absorb}.
\item $|S| = 1$, say $S = \{i\}$: the vertex $i+2$ lies outside $S$ and its
      successor $i+3$ is not $i$, so $i + 3 \notin S$; \eqref{eq:absorb} fails at
      $v = i+2$.
\item $|S| = 2$, say $S = \{i, i+2\}$: the vertex $i + 3$ lies outside $S$ and
      its successor $i+4$ is not $i$ or $i+2$, so $i+4 \notin S$;
      \eqref{eq:absorb} fails at $v = i + 3$.
\item $|S| \ge 3$: impossible, since the independence number is $2$.
\end{itemize}

Every candidate fails, so no subset $S \subseteq V$ is both independent and
out-stable. Hence $\Cfive$ has no kernel.
\end{proof}

\begin{proof}[Independent check by exhaustive enumeration]
There are $2^5 = 32$ subsets of $V$. Exactly $11$ of them are independent:
$\emptyset$; the five singletons $\{0\}, \{1\}, \{2\}, \{3\}, \{4\}$; and the
five non-adjacent pairs $\{0,2\}, \{1,3\}, \{2,4\}, \{0,3\}, \{1,4\}$. For each
one, the table below exhibits a vertex that lies outside $S$ and has no arc into
$S$, so none is out-stable; the remaining $21$ subsets are already not
independent.

\begin{center}
\begin{tabular}{c|c}
independent set $S$ & a vertex outside $S$ with no arc into $S$ \\
\hline
$\emptyset$      & $0$ (its arc goes to $1 \notin S$) \\
$\{0\}$          & $2$ (its arc goes to $3 \notin S$) \\
$\{1\}$          & $3$ (its arc goes to $4 \notin S$) \\
$\{2\}$          & $4$ (its arc goes to $0 \notin S$) \\
$\{3\}$          & $0$ (its arc goes to $1 \notin S$) \\
$\{4\}$          & $1$ (its arc goes to $2 \notin S$) \\
$\{0,2\}$        & $3$ (its arc goes to $4 \notin S$) \\
$\{1,3\}$        & $4$ (its arc goes to $0 \notin S$) \\
$\{2,4\}$        & $0$ (its arc goes to $1 \notin S$) \\
$\{0,3\}$        & $1$ (its arc goes to $2 \notin S$) \\
$\{1,4\}$        & $2$ (its arc goes to $3 \notin S$) \\
\end{tabular}
\end{center}

Thus a brute force over all $32$ subsets finds $0$ kernels, confirming
Theorem~\ref{thm:no-kernel}. The computation is reproduced by
\texttt{reproduce.py} and formalised in \texttt{lean4/Main.lean}.
\end{proof}

The case analysis above is also the general phenomenon, which we record next.

\section{The general odd-cycle obstruction}

\begin{theorem}\label{thm:cycles}
For $n \ge 3$, the directed $n$-cycle $C_n$ has a kernel if and only if $n$ is
even. When $n$ is even it has exactly two kernels, the even-indexed and the
odd-indexed vertices.
\end{theorem}

\begin{proof}
Write the vertices as $\mathbb{Z}/n$ with arcs $i \to i+1$. Independence forbids
two cyclically consecutive vertices from both lying in $S$; out-stability
\eqref{eq:absorb} forbids two cyclically consecutive vertices from both lying
outside $S$. Hence, for each $i$, exactly one of $i$ and $i+1$ lies in $S$.
Starting from $i = 0$, this forces
\[
i \in S \iff i \equiv \varepsilon \pmod 2
\]
for a fixed $\varepsilon \in \{0, 1\}$. Consistency around the cycle requires
$n \equiv 0 \pmod 2$, and then the two choices of $\varepsilon$ give the two
stated kernels. For odd $n$ no set $S$ satisfies both conditions, so $C_n$ has
no kernel.
\end{proof}

\begin{corollary}\label{cor:odd}
$C_5$, $C_7$, and $C_9$ have no kernels; in particular the conjecture fails for
every odd $n \ge 3$, with $n = 5$ the smallest value compatible with the
conjectured girth bound.
\end{corollary}

\begin{remark}[Richardson's theorem]
Richardson's theorem states that every digraph with no odd directed cycle has a
kernel. The conjectured condition ``girth at least five'' does not imply the
absence of odd directed cycles: the directed $5$-cycle is itself an odd directed
cycle and has girth $5$. The condition is therefore not a valid strengthening of
Richardson's theorem, and $\Cfive$ shows this directly.
\end{remark}

\section{Robustness: blow-ups, and either reading of girth}

Fix $m \ge 1$. The \emph{$m$-blow-up} of $\Cfive$ is obtained by replacing each
vertex $i$ by an independent part $P_i = \{(i,k) : 0 \le k < m\}$ of size $m$ and
putting an arc from every vertex of $P_i$ to every vertex of $P_{i+1 \bmod 5}$.
For $m = 1$ this is exactly $\Cfive$.

\begin{theorem}\label{thm:blowup}
For every $m \ge 1$, the $m$-blow-up of $\Cfive$ has out-degree $m$, directed
girth $5$, and no kernel.
\end{theorem}

\begin{proof}
Every vertex of a part $P_i$ has arcs to all $m$ vertices of $P_{i+1}$, and to no
other part, so the out-degree is exactly $m$. Every directed cycle projects to a
closed walk in $\Cfive$ that advances the part index by $1$ at each step, so its
length is a positive multiple of $5$; the cycle through $P_0, \dots, P_4$ has
length $5$. Hence the directed girth is $5$.

For the kernel, write $S_i = S \cap P_i$ and let $b_i = 1$ if $S_i \ne
\emptyset$ and $b_i = 0$ otherwise. If $b_i = b_{i+1} = 1$ then picking $u \in
S_i$, $v \in S_{i+1}$ gives the arc $(u, v)$, contradicting independence; so no
two consecutive $b$'s are both $1$. If $b_{i+1} = 0$ then, whenever some vertex
of $P_i$ lies outside $S$ (that is, $S_i \ne P_i$), that vertex has all its
out-arcs landing in $P_{i+1}$ and hence has no arc into $S$, contradicting
out-stability. Consequently
\begin{equation}\label{eq:blow}
b_{i+1} = 0 \implies S_i = P_i .
\end{equation}
If all $b_i = 0$, then $S_0 = \emptyset \ne P_0$ while $b_1 = 0$, contradicting
\eqref{eq:blow}. So some $b_i = 1$. Independence then forces $b_{i-1} = b_{i+1}
= 0$. Apply \eqref{eq:blow} at $i+1$: $b_{i+1} = 0$ gives $S_{i+1} = P_{i+1}
\neq \emptyset$, so $b_{i+2} = 1$. Continuing around the cycle, the $b$'s
alternate $1, 0, 1, 0, \dots$, which is consistent only if the number of parts
is even; but there are $5$ parts. Equivalently, $b_{i+5} = b_i$ would be forced
to be both $1$ and $0$. Hence no $S$ is both independent and out-stable.
\end{proof}

\begin{remark}[robustness]
Theorem~\ref{thm:blowup} shows the failure is not an artefact of out-degree $1$:
every out-degree bound $m$ admits a kernel-free digraph of girth $5$. Since the
underlying undirected graph of $\Cfive$ is a $5$-cycle whose girth is $5$ as
well, the counterexample also refutes the conjecture under the undirected
reading of ``girth''. Finally, the failure is unconditional in the sense that it
uses no asymptotic or numerical coincidence: it is a structural parity
obstruction.
\end{remark}

\section{Reproducibility}

\texttt{reproduce.py} (pure Python~3, standard library only) brute-forces all
$2^n$ subsets of the directed $n$-cycle for $n = 3, 5, 7, 9$, confirms that the
number of kernels is $0$ in every case, verifies that $\Cfive$ has out-degree
$1$ and directed and undirected girth $5$, and repeats the check for the
blow-ups of $\Cfive$ with part sizes $m = 1, 2, 3$. It prints \texttt{PASS} and
exits with status $0$ when every fact verifies, and exits non-zero otherwise.

The document is built with \texttt{tectonic main.tex}; the PDF is
\texttt{build/main.pdf}.

\section{Formalisation}

The refutation is formalised in Lean~4 in \texttt{lean4/}, using core Lean only
(\texttt{import Std}, no Mathlib, no \texttt{sorry}). The digraph is the
decidable arc predicate \texttt{arc5 i j := decide (j = i + 1)} on
\texttt{Fin 5}; subsets are indicator functions \texttt{Fin 5 → Bool}; and the
predicates \texttt{independent5}, \texttt{outStable5}, \texttt{kernel5} are
\texttt{Bool}-valued. The main theorem
\[
\texttt{no\_kernel5} : \forall S : \texttt{Fin 5} \to \texttt{Bool},\;
\texttt{kernel5 } S = \texttt{false}
\]
is proved by case analysis on the five values of $S$ --- an exhaustive split of
the $32$ subsets --- with each closed case discharged by \texttt{decide}. The
out-degree and girth facts are the theorems \texttt{outdeg5\_le\_one},
\texttt{girth5\_ge\_five}, and \texttt{girth5\_le\_five}; the explicit
$32$-subset brute force is \texttt{kernelMasks5\_nil}; and the odd cycles
$C_3$, $C_7$ are covered by \texttt{no\_kernel3}, \texttt{no\_kernel7}. The
combined statement \texttt{conjecture\_00000003485\_false} asserts that the
hypotheses of the conjecture hold for $\Cfive$ while no kernel exists.

\section{Status against the submission rules}

\begin{itemize}
\item \textbf{LaTeX source} --- present (\texttt{main.tex}), a standalone
      \texttt{article} using only \texttt{geometry}, \texttt{amsmath},
      \texttt{amssymb}, \texttt{amsthm}, \texttt{array}, and \texttt{parskip},
      designed to compile with \texttt{tectonic main.tex}.
\item \textbf{PDF document} --- \texttt{build/main.pdf}, produced by
      \texttt{tectonic main.tex}.
\item \textbf{Lean 4 project} --- present under \texttt{lean4/}, with
      \texttt{lean-toolchain} (pinned to \texttt{leanprover/lean4:v4.33.1}),
      \texttt{lakefile.toml} (library \texttt{Main}, name \texttt{tlmc3485}),
      \texttt{Main.lean}, \texttt{Check.lean}, and \texttt{README.md}. It
      contains no \texttt{sorry} and no Mathlib dependency; \texttt{Check.lean}
      reports the axioms of every theorem, and none reports \texttt{sorryAx}.
\item \textbf{Reproduction} --- \texttt{reproduce.py} is dependency-free and
      verifies all claims.
\end{itemize}

\end{document}
